\chapter{Stratégie de test}

\section{Niveaux de vérification}

\begin{table}[H]
\centering
\caption{Niveaux de test retenus}
\begin{tabularx}{\textwidth}{|L{3.2cm}|Y|L{3.4cm}|}
\hline
\rowcolor{ctmint}\textbf{Niveau} & \textbf{Objectif} & \textbf{Outils principaux} \\ \hline
Unitaire & Vérifier une règle isolée sans infrastructure réelle. & PHPUnit, pytest, Node Test Runner \\ \hline
Intégration API & Vérifier validation, persistance, autorisation, événements et isolation multi-tenant. & Laravel Test, SQLite en mémoire \\ \hline
Contrat externe & Vérifier les échanges OCPP, HMAC, paiement HTTP, webhooks et erreurs fournisseur. & pytest, WireMock, tests Laravel \\ \hline
Qualité frontend & Détecter les erreurs TypeScript, imports invalides et défauts statiques. & Oxlint, TypeScript, Vite \\ \hline
Système simulé & Vérifier un parcours vertical borne--OCPP--session--paiement. & Simulateur SAP, gateway Python, Laravel, WireMock \\ \hline
Recette manuelle & Valider l'ergonomie, les transitions visibles et les workflows par rôle. & Navigateur desktop/mobile \\ \hline
Non fonctionnel & Mesurer performance, sécurité, résilience et compatibilité. & À compléter avant production \\ \hline
\end{tabularx}
\end{table}

\section{Principes}

\begin{itemize}[leftmargin=6mm]
  \item Les autorisations sont vérifiées côté API, même si l'action est masquée
        dans l'interface.
  \item Chaque scénario multi-organisation contient au moins une tentative
        d'accès transversal.
  \item Les événements OCPP et webhooks sont rejoués pour vérifier
        l'idempotence.
  \item Les paiements restent simulés ; aucune carte ni donnée financière réelle
        n'est utilisée.
  \item Les jeux de données sont synthétiques et aucun secret n'est intégré au
        document.
  \item Un scénario manuel n'est jamais déclaré réussi sans preuve d'exécution.
\end{itemize}

\section{Priorité et sévérité}

\begin{table}[H]
\centering
\caption{Classification des anomalies}
\begin{tabularx}{\textwidth}{|C{1.7cm}|L{3.1cm}|Y|L{3cm}|}
\hline
\rowcolor{ctmint}\textbf{Niveau} & \textbf{Qualification} & \textbf{Exemple ChargeTrackr} & \textbf{Décision} \\ \hline
S1 & Bloquante & Fuite inter-organisation, paiement capturé sans session, commande envoyée à la mauvaise borne. & Recette refusée \\ \hline
S2 & Majeure & Impossible de démarrer une recharge ou de traiter une alerte critique. & Correction obligatoire \\ \hline
S3 & Mineure & Filtre incorrect, export partiel, défaut visuel gênant sans perte de données. & Correction planifiée \\ \hline
S4 & Cosmétique & Alignement, libellé ou espacement sans impact métier. & Peut être différée \\ \hline
\end{tabularx}
\end{table}

\section{Critères d'entrée et de sortie}

\begin{table}[H]
\centering
\caption{Critères de décision}
\begin{tabularx}{\textwidth}{|L{3cm}|Y|}
\hline
\rowcolor{ctmint}\textbf{Étape} & \textbf{Critères} \\ \hline
Entrée en test & Migrations applicables, variables locales définies, données synthétiques disponibles, services requis identifiés et build frontend valide. \\ \hline
Acceptation technique & Toutes les suites automatiques réussissent, aucune vulnérabilité Composer connue, aucune anomalie S1 ou S2 ouverte. \\ \hline
Acceptation fonctionnelle & Tous les cas de recette critiques et majeurs sont exécutés et réussis ou font l'objet d'une dérogation écrite. \\ \hline
Acceptation de production & Tests de charge, TLS, sauvegarde/restauration, supervision, prestataire de paiement réel et borne physique validés selon le périmètre retenu. \\ \hline
\end{tabularx}
\end{table}
